% Vietnamese translation for OI Public Library
% Source-File: IOI中国国家候选队论文/2005/李羽修/李羽修.ppt
\documentclass[11pt]{article}
\usepackage{oipl-vn}

\title{Bản trình chiếu: Thiết kế và tối ưu hàm băm}
\author{Li Yuxiu}
\date{2005}

\begin{document}
\maketitle

\origtitle{Slide companion for hash-function design}
\sourcepath{IOI China candidate team papers / 2005 / Li Yuxiu / Li Yuxiu.ppt}

\section{Phạm vi bản dịch}

Bản trình chiếu đi cùng bài viết về hàm băm. Các mốc đọc được gồm phép chia
\(h(k)=k\bmod M\), ví dụ \(k=100,M=12\), phương pháp nhân với
\(A=0.61803\), và các tên MD5, SHA1.

\section{Các cách băm cơ bản}

Với phương pháp chia, giá trị băm là \(h(k)=k\bmod M\). Cách này đơn giản,
nhưng chất lượng phụ thuộc mạnh vào việc chọn \(M\). Với phương pháp lấy bit,
khi \(M=2^r\), ta dùng \(r\) bit của khóa; ví dụ \(k=100=(1100100)_2\).

\section{Phương pháp nhân}

Một kỹ thuật khác là lấy phần lẻ của \(kA\), rồi nhân với \(M\):
\[
  h(k)=\lfloor M\{kA\}\rfloor .
\]
Slide dùng \(A=0.61803\), gần với tỉ lệ vàng, để minh họa cách phân tán khóa.

\section{Chuỗi và băm an toàn}

Với chuỗi như \texttt{IOI2005}, khóa có thể được xem như số trong hệ cơ số
hoặc dãy byte. Các hàm như MD5 và SHA1 thuộc nhóm băm mật mã; trong thi lập
trình, ta thường dùng phiên bản nhẹ hơn nhưng vẫn phải chú ý va chạm.

\end{document}
